%%%%%%%%%%%%%%%%%%%%%%%%%%%%%%%%%%%%%%%%%%%%%%%%%%%%%%%%%%%%%
%% THÉORIE DES CATÉGORIES — Chapitres 9–11 + Fin
%% Catégories monoïdales / Topos / ∞-Catégories / Appendice
%%%%%%%%%%%%%%%%%%%%%%%%%%%%%%%%%%%%%%%%%%%%%%%%%%%%%%%%%%%%%

%%%%%%%%%%%%%%%%%%%%%%%%%%%%%%%%%%%%%%%%%%%%%%%%%%%%%%%%%%%%
%%       CHAPITRE 9 : CATÉGORIES MONOÏDALES ET TRESSÉES   %%
%%%%%%%%%%%%%%%%%%%%%%%%%%%%%%%%%%%%%%%%%%%%%%%%%%%%%%%%%%%%

\chapter{Catégories monoïdales et tressées}\label{ch:monoidales}
\minitoc

\vspace{1em}

Les catégories monoïdales formalisent la notion de \og produit tensoriel\fg{}
dans un cadre catégorique général.  Elles interviennent en algèbre
(modules sur un anneau), en topologie (espaces de lacets),
en informatique théorique (types linéaires, sémantique catégorique)
et en physique mathématique (théories topologiques des champs).
Ce chapitre introduit les structures monoïdales, symétriques et tressées,
puis développe les notions de catégorie enrichie et de catégorie
cartésienne fermée, avec l'application fondamentale à la correspondance
de Curry--Howard--Lambek.

\index{catégorie monoïdale|(}

%% ============================================================
\section{Catégories monoïdales}\label{sec:cat-monoidale}
%% ============================================================

\begin{definition}[Catégorie monoïdale]\label{def:cat-monoidale}
\index{catégorie monoïdale!définition}
\index{produit tensoriel}
\index{objet unité}
\index{associateur}
\index{uniteur}
Une \textbf{catégorie monoïdale} est un sextuplet
$(\mathcal{C}, \otimes, I, \alpha, \lambda, \rho)$ où :
\begin{enumerate}[label=(\roman*)]
\item $\mathcal{C}$ est une catégorie ;
\item $\otimes \colon \mathcal{C} \times \mathcal{C} \to \mathcal{C}$
      est un bifoncteur, appelé \textbf{produit tensoriel} ;
\item $I \in \Ob(\mathcal{C})$ est l'\textbf{objet unité} ;
\item $\alpha_{A,B,C} \colon (A \otimes B) \otimes C
      \xrightarrow{\;\sim\;} A \otimes (B \otimes C)$
      est un isomorphisme naturel en $A, B, C$,
      appelé l'\textbf{associateur} ;
\item $\lambda_A \colon I \otimes A \xrightarrow{\;\sim\;} A$
      est un isomorphisme naturel en $A$,
      appelé l'\textbf{uniteur à gauche} ;
\item $\rho_A \colon A \otimes I \xrightarrow{\;\sim\;} A$
      est un isomorphisme naturel en $A$,
      appelé l'\textbf{uniteur à droite} ;
\end{enumerate}
soumis aux conditions de cohérence suivantes.

\medskip
\noindent\textbf{Diagramme du pentagone.}\index{pentagone (axiome du)}
Pour tous $A, B, C, D \in \Ob(\mathcal{C})$ :
\[
\begin{tikzcd}[column sep=small]
& (A \otimes B) \otimes (C \otimes D)
  \arrow[dr, "\alpha_{A,B,C \otimes D}"] & \\
((A \otimes B) \otimes C) \otimes D
  \arrow[ur, "\alpha_{A \otimes B,C,D}"]
  \arrow[d, "\alpha_{A,B,C} \otimes \id_D"'] & &
A \otimes (B \otimes (C \otimes D)) \\
(A \otimes (B \otimes C)) \otimes D
  \arrow[rr, "\alpha_{A,B \otimes C,D}"'] & &
A \otimes ((B \otimes C) \otimes D)
  \arrow[u, "\id_A \otimes \alpha_{B,C,D}"']
\end{tikzcd}
\]

\noindent\textbf{Diagramme du triangle.}\index{triangle (axiome du)}
Pour tous $A, B \in \Ob(\mathcal{C})$ :
\[
\begin{tikzcd}
(A \otimes I) \otimes B
  \arrow[rr, "\alpha_{A,I,B}"]
  \arrow[dr, "\rho_A \otimes \id_B"'] & &
A \otimes (I \otimes B)
  \arrow[dl, "\id_A \otimes \lambda_B"] \\
& A \otimes B &
\end{tikzcd}
\]
\end{definition}

\begin{definition}[Catégorie monoïdale stricte]\label{def:monoidale-stricte}
\index{catégorie monoïdale!stricte}
Une catégorie monoïdale est dite \textbf{stricte} si les isomorphismes
$\alpha$, $\lambda$ et $\rho$ sont tous des identités, c'est-à-dire
si $(A \otimes B) \otimes C = A \otimes (B \otimes C)$,
$I \otimes A = A$ et $A \otimes I = A$ pour tous objets $A, B, C$.
\end{definition}

\begin{theoreme}[Cohérence de Mac Lane]\label{thm:coherence-maclane}
\index{théorème de cohérence!Mac Lane}
\index{Mac Lane!théorème de cohérence}
Soit $(\mathcal{C}, \otimes, I, \alpha, \lambda, \rho)$ une catégorie
monoïdale.
\begin{enumerate}[label=(\alph*)]
\item Tout diagramme dont les flèches sont des composées d'instances de
      $\alpha$, $\lambda$, $\rho$ et de leurs inverses commute.
\item Toute catégorie monoïdale est monoïdalement équivalente à une
      catégorie monoïdale stricte.
\end{enumerate}
\end{theoreme}

\begin{proof}
La partie (a) est le théorème original de Mac Lane (1963).
On montre d'abord que tout diagramme de cohérence se réduit à un diagramme
dans la catégorie libre monoïdale sur un générateur, puis on vérifie
la commutativité dans ce cas par induction sur la complexité des
expressions.  La partie (b) se déduit en construisant une catégorie
monoïdale stricte $\mathcal{C}_s$ avec une équivalence monoïdale
$\mathcal{C} \simeq \mathcal{C}_s$ (cf.\ Joyal--Street).
\end{proof}

\begin{remarque}\label{rem:coherence-pratique}
\index{théorème de cohérence!conséquence pratique}
En pratique, le théorème de cohérence nous autorise à omettre les
parenthèses et à identifier $I \otimes A$, $A \otimes I$ et $A$
sans ambiguïté.  On écrit simplement $A \otimes B \otimes C$
au lieu de $(A \otimes B) \otimes C$ ou $A \otimes (B \otimes C)$.
\end{remarque}

%% ============================================================
\section{Exemples fondamentaux}\label{sec:monoidale-exemples}
%% ============================================================

\begin{exemple}[Catégories monoïdales classiques]\label{ex:monoidales-classiques}
\index{catégorie monoïdale!exemples}
\begin{enumerate}[label=(\roman*)]
\item \textbf{$(\Set, \times, \{*\})$} :\index{Set@$\Set$!monoïdale}
      la catégorie des ensembles avec le produit cartésien.
      L'objet unité est un singleton.

\item \textbf{$(\Ab, \otimes_\Z, \Z)$} :\index{Ab@$\Ab$!monoïdale}
      la catégorie des groupes abéliens avec le produit tensoriel
      sur $\Z$.

\item \textbf{$(R\text{-}\Mod, \otimes_R, R)$} :\index{Mod@$\Mod$!monoïdale}
      pour un anneau commutatif $R$, la catégorie des $R$-modules
      avec le produit tensoriel.

\item \textbf{$(\Vect_\K, \otimes_\K, \K)$} :\index{Vect@$\Vect$!monoïdale}
      la catégorie des $\K$-espaces vectoriels.

\item \textbf{$(\Cat, \times, \mathbf{1})$} :\index{Cat@$\Cat$!monoïdale}
      la catégorie des petites catégories avec le produit.

\item \textbf{$(\Set, \sqcup, \varnothing)$} :
      les ensembles avec l'union disjointe.  L'objet unité
      est l'ensemble vide.

\item \textbf{$([\mathcal{C}, \mathcal{C}], \circ, \id_{\mathcal{C}})$} :
      \index{catégorie d'endofoncteurs}
      la catégorie des endofoncteurs de $\mathcal{C}$ avec la
      composition.  Cette catégorie monoïdale est stricte.
\end{enumerate}
\end{exemple}

\begin{exemple}[Structures monoïdales sur les espaces vectoriels gradués]
\label{ex:ev-gradues-monoidal}
\index{espace vectoriel gradué!structure monoïdale}
Soit $\Vect_\K^{\Z}$ la catégorie des espaces vectoriels $\Z$-gradués.
Le produit tensoriel gradué est défini par
\[
  (V \otimes W)_n = \bigoplus_{p+q=n} V_p \otimes_\K W_q.
\]
L'objet unité est $\K$ concentré en degré $0$.
Cette structure monoïdale joue un rôle central en algèbre homologique
et en topologie algébrique.
\end{exemple}

%% ============================================================
\section{Catégories monoïdales symétriques et tressées}
\label{sec:symetrique-tressee}
%% ============================================================

\begin{definition}[Catégorie monoïdale tressée]\label{def:monoidale-tressee}
\index{catégorie monoïdale!tressée}
\index{tressage}
Une \textbf{catégorie monoïdale tressée} est une catégorie monoïdale
$(\mathcal{C}, \otimes, I)$ munie d'un isomorphisme naturel
\[
  \sigma_{A,B} \colon A \otimes B \xrightarrow{\;\sim\;} B \otimes A,
\]
appelé \textbf{tressage}, satisfaisant les deux diagrammes hexagonaux
suivants.

\medskip
\noindent\textbf{Premier hexagone} :\index{hexagone (axiome de l')}
\[
\begin{tikzcd}[column sep=1.8em]
(A \otimes B) \otimes C
  \arrow[r, "\sigma_{A,B} \otimes \id"]
  \arrow[d, "\alpha"'] &
(B \otimes A) \otimes C
  \arrow[r, "\alpha"] &
B \otimes (A \otimes C)
  \arrow[d, "\id \otimes \sigma_{A,C}"] \\
A \otimes (B \otimes C)
  \arrow[r, "\sigma_{A,B \otimes C}"'] &
(B \otimes C) \otimes A
  \arrow[r, "\alpha"'] &
B \otimes (C \otimes A)
\end{tikzcd}
\]

\noindent\textbf{Second hexagone} : le diagramme analogue avec
$\sigma^{-1}$ et les associateurs inversés.
\end{definition}

\begin{definition}[Catégorie monoïdale symétrique]\label{def:monoidale-symetrique}
\index{catégorie monoïdale!symétrique}
\index{symétrie (tressage)}
Une catégorie monoïdale tressée est dite \textbf{symétrique} si
$\sigma_{B,A} \circ \sigma_{A,B} = \id_{A \otimes B}$
pour tous $A, B$.
\end{definition}

\begin{exemple}\label{ex:symetrique-vs-tressee}
\begin{enumerate}[label=(\roman*)]
\item $(\Set, \times, \{*\})$, $(\Ab, \otimes_\Z, \Z)$,
      $(\Vect_\K, \otimes_\K, \K)$ sont symétriques.
\item La catégorie des représentations d'un groupe quantique
      $U_q(\mathfrak{g})$ pour $q$ générique est tressée mais
      \emph{pas} symétrique en général.
      \index{groupe quantique}
\end{enumerate}
\end{exemple}

%% ============================================================
\section{Objets monoïdes}\label{sec:objets-monoides}
%% ============================================================

\begin{definition}[Monoïde dans une catégorie monoïdale]
\label{def:objet-monoide}
\index{monoïde!dans une catégorie monoïdale}
\index{objet monoïde}
Soit $(\mathcal{C}, \otimes, I)$ une catégorie monoïdale.
Un \textbf{objet monoïde} (ou \textbf{monoïde dans $\mathcal{C}$})
est un triplet $(M, \mu, \eta)$ où :
\begin{enumerate}[label=(\roman*)]
\item $M \in \Ob(\mathcal{C})$ ;
\item $\mu \colon M \otimes M \to M$ est la \textbf{multiplication} ;
\item $\eta \colon I \to M$ est l'\textbf{unité} ;
\end{enumerate}
satisfaisant les axiomes d'associativité et d'unité :
\[
\begin{tikzcd}
(M \otimes M) \otimes M
  \arrow[r, "\alpha"]
  \arrow[d, "\mu \otimes \id"'] &
M \otimes (M \otimes M)
  \arrow[d, "\id \otimes \mu"] \\
M \otimes M \arrow[r, "\mu"'] &
M
\end{tikzcd}
\qquad
\begin{tikzcd}
I \otimes M
  \arrow[r, "\eta \otimes \id"]
  \arrow[dr, "\lambda"'] &
M \otimes M
  \arrow[d, "\mu"] &
M \otimes I
  \arrow[l, "\id \otimes \eta"']
  \arrow[dl, "\rho"] \\
& M &
\end{tikzcd}
\]
\end{definition}

\begin{exemple}[Monoïdes dans diverses catégories]\label{ex:objets-monoides}
\index{objet monoïde!exemples}
\begin{enumerate}[label=(\roman*)]
\item Un monoïde dans $(\Set, \times, \{*\})$ est un monoïde
      au sens usuel.
\item Un monoïde dans $(\Ab, \otimes_\Z, \Z)$ est un anneau.
      \index{anneau!comme monoïde dans Ab}
\item Un monoïde dans $(R\text{-}\Mod, \otimes_R, R)$ est une
      $R$-algèbre.\index{algèbre!comme monoïde}
\item Un monoïde dans
      $([\mathcal{C}, \mathcal{C}], \circ, \id_\mathcal{C})$
      est une monade.\index{monade!comme monoïde}
\end{enumerate}
\end{exemple}

\begin{definition}[Comonoïde]\label{def:comonoide}
\index{comonoïde}
Dualement, un \textbf{comonoïde} dans $(\mathcal{C}, \otimes, I)$
est un triplet $(C, \Delta, \varepsilon)$ avec
$\Delta \colon C \to C \otimes C$ (comultiplication) et
$\varepsilon \colon C \to I$ (coünité), satisfaisant les axiomes duaux.
Un comonoïde dans $(\Vect_\K, \otimes_\K, \K)$ est une coalgèbre.
\index{coalgèbre}
\end{definition}

%% ============================================================
\section{Catégories enrichies}\label{sec:categories-enrichies}
%% ============================================================

\begin{definition}[Catégorie enrichie]\label{def:cat-enrichie}
\index{catégorie enrichie}
\index{catégorie-$\mathcal{V}$}
Soit $(\mathcal{V}, \otimes, I)$ une catégorie monoïdale.
Une \textbf{catégorie enrichie sur $\mathcal{V}$} (ou
\textbf{$\mathcal{V}$-catégorie}) $\mathcal{C}$ consiste en :
\begin{enumerate}[label=(\roman*)]
\item une collection d'objets $\Ob(\mathcal{C})$ ;
\item pour chaque paire $(A, B)$, un objet
      $\mathcal{C}(A,B) \in \Ob(\mathcal{V})$
      (l'\textbf{objet de morphismes}) ;
\item pour chaque triplet $(A, B, C)$, un morphisme de composition
      $\circ_{A,B,C} \colon \mathcal{C}(B,C) \otimes \mathcal{C}(A,B)
      \to \mathcal{C}(A,C)$ dans $\mathcal{V}$ ;
\item pour chaque objet $A$, un morphisme identité
      $j_A \colon I \to \mathcal{C}(A,A)$ dans $\mathcal{V}$ ;
\end{enumerate}
satisfaisant les axiomes d'associativité et d'unité
(exprimés comme commutativité de diagrammes dans $\mathcal{V}$).
\end{definition}

\begin{exemple}[Enrichissements classiques]\label{ex:enrichissements}
\index{catégorie enrichie!exemples}
\begin{enumerate}[label=(\roman*)]
\item $\Set$-catégories $=$ catégories ordinaires
      (localement petites).
\item $\Ab$-catégories $=$ catégories préadditives.
      \index{catégorie préadditive}
\item $\Vect_\K$-catégories $=$ catégories $\K$-linéaires.
\item $\Cat$-catégories $=$ $2$-catégories (strictes).
      \index{2-catégorie}
\item $([0,\infty], \geq, +, 0)$-catégories $=$ espaces métriques
      généralisés (au sens de Lawvere).
      \index{espace métrique!de Lawvere}
\end{enumerate}
\end{exemple}

%% ============================================================
\section{Catégories fermées}\label{sec:categories-fermees}
%% ============================================================

\begin{definition}[Catégorie monoïdale fermée]\label{def:monoidale-fermee}
\index{catégorie monoïdale!fermée}
\index{Hom interne}
\index{exponentiation}
Une catégorie monoïdale $(\mathcal{C}, \otimes, I)$ est dite
\textbf{fermée (à droite)} si pour chaque objet $B$,
le foncteur $- \otimes B \colon \mathcal{C} \to \mathcal{C}$
admet un adjoint à droite, noté $[B, -]$ ou $\underline{\Hom}(B, -)$ :
\[
  \Hom_\mathcal{C}(A \otimes B, C)
  \;\cong\;
  \Hom_\mathcal{C}(A, [B, C])
\]
naturellement en $A$ et $C$.
L'objet $[B, C]$ est appelé le \textbf{Hom interne}.
\end{definition}

\begin{definition}[Catégorie cartésienne fermée]\label{def:cartesienne-fermee}
\index{catégorie cartésienne fermée}
\index{CCC|see{catégorie cartésienne fermée}}
Une catégorie est \textbf{cartésienne fermée} (CCC) si :
\begin{enumerate}[label=(\roman*)]
\item elle possède un objet terminal $1$ ;
\item elle possède tous les produits binaires $A \times B$ ;
\item elle est monoïdale fermée pour la structure
      $(\mathcal{C}, \times, 1)$, c'est-à-dire les exponentielles
      $B^A = [A, B]$ existent :
      \[
        \Hom(A \times B, C) \cong \Hom(A, C^B).
      \]
\end{enumerate}
\end{definition}

\begin{exemple}[Catégories cartésiennes fermées]\label{ex:ccc}
\index{catégorie cartésienne fermée!exemples}
\begin{enumerate}[label=(\roman*)]
\item $\Set$ avec $B^A = \Set(A, B)$.
\item $\Cat$ avec $\mathcal{D}^\mathcal{C} = \Fun(\mathcal{C}, \mathcal{D})$.
\item La catégorie des graphes orientés.
\item La catégorie des ensembles simpliciaux
      $\mathbf{sSet} = \Fun(\Delta^\op, \Set)$.
      \index{ensemble simplicial}
\item Tout topos élémentaire (cf.\ chapitre~\ref{ch:topos}).
\end{enumerate}
\end{exemple}

\begin{proposition}[Évaluation et curryfication]\label{prop:eval-curry}
\index{évaluation (morphisme d')}
\index{curryfication}
Dans une catégorie cartésienne fermée, on dispose de :
\begin{enumerate}[label=(\alph*)]
\item un morphisme d'\textbf{évaluation}
      $\mathrm{ev}_{A,B} \colon B^A \times A \to B$ ;
\item pour tout $f \colon A \times B \to C$, un unique
      $\tilde{f} \colon A \to C^B$ (le \textbf{curryfié} de $f$)
      tel que $\mathrm{ev} \circ (\tilde{f} \times \id_B) = f$.
\end{enumerate}
\end{proposition}

\begin{proof}
Le morphisme d'évaluation est la coünité de l'adjonction
$(-) \times A \dashv (-)^A$.  La curryfication est le passage au
transposé adjoint.
\end{proof}

%% ============================================================
\section{Correspondance de Curry--Howard--Lambek}\label{sec:curry-howard-lambek}
%% ============================================================

\index{correspondance de Curry--Howard--Lambek|(}
\index{Curry--Howard--Lambek|see{correspondance de Curry--Howard--Lambek}}

La correspondance de Curry--Howard--Lambek établit un dictionnaire
entre trois domaines :

\medskip
\begin{center}
\renewcommand{\arraystretch}{1.3}
\begin{tabular}{lll}
\toprule
\textbf{Logique} & \textbf{Typage} & \textbf{Catégories} \\
\midrule
Proposition $A$ & Type $A$ & Objet $A$ \\
Preuve de $A$ & Terme de type $A$ & Morphisme $1 \to A$ \\
$A \Rightarrow B$ & Type fonction $A \to B$ & Exponentielle $B^A$ \\
$A \wedge B$ & Type produit $A \times B$ & Produit $A \times B$ \\
Vrai ($\top$) & Type unité & Objet terminal $1$ \\
Modus ponens & Application & Évaluation \\
Introduction de $\Rightarrow$ & $\lambda$-abstraction & Curryfication \\
\bottomrule
\end{tabular}
\end{center}

\begin{theoreme}[Lambek]\label{thm:lambek}
\index{Lambek!théorème de}
Il y a une équivalence de catégories entre :
\begin{enumerate}[label=(\alph*)]
\item les théories de type du $\lambda$-calcul simplement typé
      (avec produits) ;
\item les catégories cartésiennes fermées.
\end{enumerate}
Plus précisément, la catégorie syntaxique d'un $\lambda$-calcul
simplement typé est cartésienne fermée, et réciproquement,
le langage interne d'une CCC est un $\lambda$-calcul simplement typé.
\index{lambda-calcul@$\lambda$-calcul!simplement typé}
\end{theoreme}

\begin{proof}[Esquisse]
On construit la catégorie syntaxique $\mathcal{C}_T$ d'une théorie $T$ :
les objets sont les types, les morphismes $A \to B$ sont les termes
$x : A \vdash t : B$ modulo $\beta\eta$-équivalence.
Le produit de types donne le produit catégorique,
et le type fonction donne l'exponentielle.
Réciproquement, à une CCC $\mathcal{C}$ on associe la théorie
dont les types sont les objets et dont les termes sont les morphismes.
\end{proof}

\index{correspondance de Curry--Howard--Lambek|)}

%% ============================================================
\section{Exercices}\label{sec:exo-monoidales}
%% ============================================================

\begin{exercice}\label{exo:monoidale-strict}
\index{catégorie monoïdale!stricte}
Montrer que la catégorie $(\mathbf{Mat}_\K, \otimes, 1)$, dont les
objets sont les entiers naturels et les morphismes $m \to n$ sont
les matrices $n \times m$ à coefficients dans $\K$, avec le produit
de Kronecker, est monoïdale stricte.
\end{exercice}

\begin{exercice}\label{exo:monoide-dans-Ab}
Vérifier en détail que les monoïdes dans $(\Ab, \otimes_\Z, \Z)$
sont exactement les anneaux (unitaires).
\end{exercice}

\begin{exercice}\label{exo:ccc-pas-Set}
Montrer que la catégorie $\Top$ des espaces topologiques n'est
\emph{pas} cartésienne fermée.
\textit{Indication} : considérer le produit $\Q \times \Q$ et la
topologie compacte-ouverte.
\index{Top@$\Top$!pas cartésienne fermée}
\end{exercice}

\begin{exercice}\label{exo:enrichie-metrique}
\index{espace métrique!de Lawvere}
Soit $\mathcal{V} = ([0, \infty], \geq)$ munie de la structure
monoïdale $(+, 0)$.  Montrer qu'une $\mathcal{V}$-catégorie est
exactement un espace quasi-métrique généralisé au sens de Lawvere :
$d(x,x) = 0$ et $d(x,z) \leq d(x,y) + d(y,z)$.
\end{exercice}

\begin{exercice}[Groupe de tresses et catégorie tressée]
\label{exo:tresses}
\index{groupe de tresses}
Soit $\mathcal{B}$ la catégorie dont les objets sont les entiers
$n \geq 0$ et dont les morphismes $n \to n$ forment le groupe de
tresses $B_n$.  Montrer que $\mathcal{B}$ est monoïdale (pour
l'addition) et tressée.
\end{exercice}

\begin{exercice}\label{exo:curry-logique}
En utilisant la correspondance de Curry--Howard--Lambek, interpréter
la tautologie $(A \Rightarrow (B \Rightarrow C)) \Leftrightarrow
((A \wedge B) \Rightarrow C)$ comme un isomorphisme dans une CCC.
\end{exercice}

\index{catégorie monoïdale|)}


%%%%%%%%%%%%%%%%%%%%%%%%%%%%%%%%%%%%%%%%%%%%%%%%%%%%%%%%%%%%
%%      CHAPITRE 10 : TOPOS DE GROTHENDIECK               %%
%%%%%%%%%%%%%%%%%%%%%%%%%%%%%%%%%%%%%%%%%%%%%%%%%%%%%%%%%%%%

\chapter{Topos de Grothendieck --- Introduction}\label{ch:topos}
\minitoc

\vspace{1em}

La théorie des topos, initiée par Grothendieck dans le cadre de la
géométrie algébrique et reformulée par Lawvere et Tierney dans un cadre
axiomatique, constitue l'un des sommets de la théorie des catégories.
Un topos unifie la géométrie (faisceaux sur un espace),
l'algèbre (faisceaux sur un site) et la logique (univers de
mathématiques intuitionnistes).  Ce chapitre présente les fondements :
sites, topologies de Grothendieck, faisceaux, topos élémentaires,
classificateur de sous-objets et logique interne.

\index{topos|(}

%% ============================================================
\section{Sites et topologies de Grothendieck}\label{sec:sites}
%% ============================================================

\begin{definition}[Crible]\label{def:crible}
\index{crible}
\index{crible!sous-foncteur}
Soit $\mathcal{C}$ une petite catégorie et $C \in \Ob(\mathcal{C})$.
Un \textbf{crible} sur $C$ est un sous-foncteur
$S \hookrightarrow \Hom_\mathcal{C}(-, C)$
du préfaisceau représentable, c'est-à-dire une famille de morphismes
de but $C$ stable par précomposition : si $f \colon D \to C$ est dans
$S$ et $g \colon E \to D$ est un morphisme quelconque, alors
$f \circ g \in S$.
\end{definition}

\begin{definition}[Topologie de Grothendieck]\label{def:topologie-grothendieck}
\index{topologie de Grothendieck}
\index{site}
Une \textbf{topologie de Grothendieck} $J$ sur une petite catégorie
$\mathcal{C}$ est la donnée, pour chaque objet $C \in \Ob(\mathcal{C})$,
d'une collection $J(C)$ de cribles sur $C$ (appelés \textbf{cribles
couvrants}), satisfaisant les axiomes :
\begin{enumerate}[label=(T\arabic*)]
\item \textbf{Maximalité} : le crible maximal
      $\Hom_\mathcal{C}(-, C)$ est dans $J(C)$.
\item \textbf{Stabilité} : si $S \in J(C)$ et $f \colon D \to C$,
      alors le crible tiré en arrière
      $f^*(S) = \{g \colon E \to D \mid f \circ g \in S\}$
      est dans $J(D)$.
\item \textbf{Transitivité} : si $S \in J(C)$ et $R$ est un crible
      sur $C$ tel que pour tout $f \colon D \to C$ dans $S$ on a
      $f^*(R) \in J(D)$, alors $R \in J(C)$.
\end{enumerate}
Un couple $(\mathcal{C}, J)$ est appelé un \textbf{site}.
\end{definition}

\begin{exemple}[Sites fondamentaux]\label{ex:sites}
\index{site!exemples}
\begin{enumerate}[label=(\roman*)]
\item \textbf{Site ouvert.}\index{site!ouvert}
      Soit $X$ un espace topologique.
      La catégorie $\mathsf{Ouv}(X)$ des ouverts de $X$
      (avec les inclusions) munie de la topologie où un crible
      sur $U$ est couvrant si et seulement si les ouverts qu'il
      contient recouvrent $U$.

\item \textbf{Topologie étale.}\index{topologie étale}
      Sur la catégorie des schémas de type fini sur un corps $k$,
      les familles couvrantes sont les familles surjectives de
      morphismes étales.

\item \textbf{Topologie fppf.}\index{topologie fppf}
      Les familles couvrantes sont les familles surjectives de
      morphismes fidèlement plats de présentation finie.

\item \textbf{Topologie de Zariski.}\index{topologie de Zariski}
      Les familles couvrantes sont les recouvrements par des
      ouverts de Zariski.

\item \textbf{Topologie grossière.}\index{topologie grossière}
      Seul le crible maximal est couvrant.  Tout préfaisceau est
      un faisceau pour cette topologie.

\item \textbf{Topologie canonique.}\index{topologie canonique}
      La plus fine topologie pour laquelle tout foncteur
      représentable est un faisceau.
\end{enumerate}
\end{exemple}

%% ============================================================
\section{Préfaisceaux et faisceaux}\label{sec:faisceaux-site}
%% ============================================================

\begin{definition}[Préfaisceau]\label{def:prefaisceau-site}
\index{préfaisceau!sur une catégorie}
Un \textbf{préfaisceau} (d'ensembles) sur une petite catégorie
$\mathcal{C}$ est un foncteur $F \colon \mathcal{C}^\op \to \Set$.
La catégorie des préfaisceaux est notée
$\widehat{\mathcal{C}} = \Fun(\mathcal{C}^\op, \Set)$.
\index{catégorie de préfaisceaux}
\end{definition}

\begin{definition}[Condition de faisceau]\label{def:faisceau-site}
\index{faisceau!sur un site}
\index{condition de recollement}
Soit $(\mathcal{C}, J)$ un site.  Un préfaisceau
$F \colon \mathcal{C}^\op \to \Set$ est un \textbf{faisceau}
(pour la topologie $J$) si pour tout objet $C$ et tout crible
couvrant $S \in J(C)$, le morphisme canonique
\[
  F(C) \longrightarrow \lim_{\substack{(f \colon D \to C) \in S}} F(D)
\]
est un isomorphisme.  Autrement dit, toute famille compatible d'éléments
indexée par $S$ se recolle de manière unique.
\end{definition}

\begin{remarque}\label{rem:faisceau-egalisateur}
\index{diagramme égalisateur}
Lorsque la topologie est engendrée par des familles couvrantes
$\{f_i \colon U_i \to C\}_{i \in I}$, la condition de faisceau
se reformule par l'exactitude du diagramme égalisateur classique :
\[
\begin{tikzcd}
F(C) \arrow[r] &
\displaystyle\prod_{i \in I} F(U_i)
  \arrow[r, shift left=0.5ex]
  \arrow[r, shift right=0.5ex] &
\displaystyle\prod_{i,j \in I} F(U_i \times_C U_j)
\end{tikzcd}
\]
\end{remarque}

%% ============================================================
\section{Topos de Grothendieck}\label{sec:topos-grothendieck}
%% ============================================================

\begin{definition}[Topos de Grothendieck]\label{def:topos-grothendieck}
\index{topos!de Grothendieck}
\index{Sh@$\mathrm{Sh}(\mathcal{C}, J)$}
Un \textbf{topos de Grothendieck} est une catégorie équivalente
à la catégorie $\mathrm{Sh}(\mathcal{C}, J)$ des faisceaux d'ensembles
sur un site $(\mathcal{C}, J)$.
\end{definition}

\begin{proposition}[Faisceautisation]\label{prop:faisceautisation}
\index{faisceautisation}
\index{foncteur faisceau associé}
L'inclusion $\iota \colon \mathrm{Sh}(\mathcal{C}, J) \hookrightarrow
\widehat{\mathcal{C}}$ admet un adjoint à gauche
$a \colon \widehat{\mathcal{C}} \to \mathrm{Sh}(\mathcal{C}, J)$,
appelé \textbf{foncteur faisceau associé} ou \textbf{faisceautisation}.
L'adjonction $a \dashv \iota$ vérifie que $a$ est exact à gauche.
\end{proposition}

\begin{proof}[Esquisse]
On construit $a$ en deux étapes d'une procédure de type «~plus~» ($+$).
Pour un préfaisceau $F$, on définit
$F^+(C) = \colim_{S \in J(C)} \Hom(S, F)$.
Le foncteur $F \mapsto (F^+)^+$ donne le faisceau associé.
\end{proof}

\begin{exemple}[Topos fondamentaux]\label{ex:topos}
\index{topos!exemples}
\begin{enumerate}[label=(\roman*)]
\item \textbf{Topos des faisceaux sur un espace.}
      \index{topos!de faisceaux sur un espace}
      Pour un espace topologique $X$,
      $\mathrm{Sh}(X) = \mathrm{Sh}(\mathsf{Ouv}(X), J_{\text{ouv}})$
      est le topos classique.

\item \textbf{Topos étale.}
      \index{topos!étale}
      Pour un schéma $X$, le petit site étale donne le topos étale
      $X_{\text{ét}}$, central en cohomologie étale.

\item \textbf{Topos des préfaisceaux.}
      \index{topos!de préfaisceaux}
      $\widehat{\mathcal{C}} = \Fun(\mathcal{C}^\op, \Set)$ est un
      topos (pour la topologie grossière).

\item \textbf{Topos classifiant.}
      \index{topos!classifiant}
      Pour un groupe $G$, la catégorie $G\text{-}\Set$ des $G$-ensembles
      est un topos.
\end{enumerate}
\end{exemple}

%% ============================================================
\section{Topos élémentaire}\label{sec:topos-elementaire}
%% ============================================================

\begin{definition}[Topos élémentaire]\label{def:topos-elementaire}
\index{topos!élémentaire}
\index{Lawvere}
\index{Tierney}
Un \textbf{topos élémentaire} est une catégorie $\mathcal{E}$ vérifiant :
\begin{enumerate}[label=(\roman*)]
\item $\mathcal{E}$ possède toutes les limites finies ;
\item $\mathcal{E}$ est cartésienne fermée ;
\item $\mathcal{E}$ possède un \textbf{classificateur de sous-objets}
      $\Omega$ (cf.\ définition~\ref{def:classificateur}).
\end{enumerate}
\end{definition}

\begin{definition}[Classificateur de sous-objets]\label{def:classificateur}
\index{classificateur de sous-objets}
\index{Omega@$\Omega$ (classificateur)}
\index{morphisme caractéristique}
Un objet $\Omega$ muni d'un monomorphisme
$\mathsf{vrai} \colon 1 \rightarrowtail \Omega$
est un \textbf{classificateur de sous-objets} si pour tout
monomorphisme $m \colon S \rightarrowtail A$, il existe un unique
morphisme $\chi_m \colon A \to \Omega$ (le \textbf{morphisme
caractéristique}) tel que le carré suivant est un tiré en arrière :
\[
\begin{tikzcd}
S \arrow[r] \arrow[d, tail, "m"'] & 1 \arrow[d, tail, "\mathsf{vrai}"] \\
A \arrow[r, "\chi_m"'] & \Omega
\end{tikzcd}
\]
\end{definition}

\begin{exemple}[Classificateur dans $\Set$]\label{ex:omega-set}
\index{classificateur de sous-objets!dans Set}
Dans $\Set$, le classificateur est $\Omega = \{0, 1\}$ avec
$\mathsf{vrai} \colon \{*\} \to \{0,1\}$ envoyant $*$ sur $1$.
Pour un sous-ensemble $S \subseteq A$, le morphisme caractéristique
est la fonction indicatrice $\chi_S$.
\end{exemple}

\begin{proposition}\label{prop:topos-groth-elementaire}
\index{topos!de Grothendieck!est élémentaire}
Tout topos de Grothendieck est un topos élémentaire.  La réciproque
est fausse en général.
\end{proposition}

\begin{proof}[Esquisse]
Un topos de Grothendieck possède toutes les limites (et colimites),
est cartésien fermé, et le foncteur $\mathrm{Sub}(-) \colon
\mathcal{E}^\op \to \Set$ est représentable par un objet $\Omega$.
Pour la non-réciproque, la catégorie des ensembles finis est un
topos élémentaire mais pas de Grothendieck (pas de colimites
quelconques).
\end{proof}

%% ============================================================
\section{Logique interne d'un topos}\label{sec:logique-interne}
%% ============================================================

\index{logique interne|(}
\index{logique intuitionniste}

La structure de topos élémentaire induit une logique interne
qui est, en général, \emph{intuitionniste} : le tiers exclu
$A \vee \neg A$ n'est pas nécessairement valide.

\begin{definition}[Algèbre de Heyting interne]\label{def:heyting-interne}
\index{algèbre de Heyting!interne}
\index{sous-objet!treillis de}
Dans un topos élémentaire $\mathcal{E}$, pour chaque objet $A$,
le treillis des sous-objets $\mathrm{Sub}(A)$ est une algèbre de Heyting.
Les opérations logiques sont :
\begin{enumerate}[label=(\roman*)]
\item $\wedge$ : intersection de sous-objets (tiré en arrière le
      long de la diagonale) ;
\item $\vee$ : union de sous-objets (image de la somme
      amalgamée) ;
\item $\Rightarrow$ : implication de Heyting (adjoint à droite de
      $\wedge$) ;
\item $\neg$ : négation définie par $\neg S = (S \Rightarrow \bot)$.
\end{enumerate}
\end{definition}

\begin{proposition}[Structure algébrique de $\Omega$]\label{prop:omega-heyting}
\index{Omega@$\Omega$!algèbre de Heyting}
Dans un topos élémentaire, l'objet $\Omega$ est un objet d'algèbre
de Heyting interne.  Les morphismes
$\wedge, \vee, \Rightarrow \colon \Omega \times \Omega \to \Omega$
et $\neg \colon \Omega \to \Omega$ satisfont les axiomes d'algèbre
de Heyting.  L'objet $\Omega$ est une algèbre de Boole interne si et
seulement si la logique du topos est classique.
\end{proposition}

\begin{remarque}\label{rem:topos-logique-classique}
\index{topos!booléen}
Un topos $\mathcal{E}$ est dit \textbf{booléen} si $\Omega$ est
une algèbre de Boole interne, c'est-à-dire si
$\neg \neg = \id_\Omega$.  Les conditions suivantes sont
équivalentes :
\begin{enumerate}[label=(\roman*)]
\item $\mathcal{E}$ est booléen ;
\item tout sous-objet admet un complément ;
\item $\mathrm{Sub}(A)$ est une algèbre de Boole pour tout $A$.
\end{enumerate}
Le topos $\Set$ est booléen.  Le topos $\mathrm{Sh}(X)$ pour un
espace $X$ est booléen si et seulement si $X$ est
extrêmement disconnecté (pour les faisceaux de Stone).
\end{remarque}

%% ============================================================
\section{Morphismes géométriques}\label{sec:morphismes-geometriques}
%% ============================================================

\begin{definition}[Morphisme géométrique]\label{def:morphisme-geometrique}
\index{morphisme géométrique}
\index{image directe}
\index{image inverse}
Un \textbf{morphisme géométrique} $f \colon \mathcal{F} \to \mathcal{E}$
entre topos est un couple de foncteurs
$(f^* \colon \mathcal{E} \to \mathcal{F},\;
  f_* \colon \mathcal{F} \to \mathcal{E})$
tel que :
\begin{enumerate}[label=(\roman*)]
\item $f^* \dashv f_*$ (adjonction) ;
\item $f^*$ est exact à gauche (préserve les limites finies).
\end{enumerate}
Le foncteur $f^*$ est l'\textbf{image inverse} et $f_*$ est
l'\textbf{image directe}.
\end{definition}

\begin{exemple}\label{ex:morphismes-geo}
\index{morphisme géométrique!exemples}
\begin{enumerate}[label=(\roman*)]
\item Pour une application continue $f \colon X \to Y$,
      l'image directe et l'image inverse des faisceaux donnent
      un morphisme géométrique $\mathrm{Sh}(X) \to \mathrm{Sh}(Y)$.
\item Le morphisme global $\Gamma \colon \mathcal{E} \to \Set$
      défini par $\Gamma = \Hom_\mathcal{E}(1, -)$ est l'image
      directe d'un morphisme géométrique
      $\gamma \colon \mathcal{E} \to \Set$
      (le \textbf{point global}).
      \index{point global}
\end{enumerate}
\end{exemple}

%% ============================================================
\section{Théorème de Giraud}\label{sec:giraud}
%% ============================================================

\begin{theoreme}[Giraud]\label{thm:giraud}
\index{théorème de Giraud}
\index{Giraud!théorème de}
Soit $\mathcal{E}$ une catégorie.  Les conditions suivantes sont
équivalentes :
\begin{enumerate}[label=(\alph*)]
\item $\mathcal{E}$ est un topos de Grothendieck (i.e.\ $\mathcal{E}
      \simeq \mathrm{Sh}(\mathcal{C}, J)$ pour un petit site
      $(\mathcal{C}, J)$) ;
\item $\mathcal{E}$ vérifie :
  \begin{enumerate}[label=(\roman*)]
  \item $\mathcal{E}$ possède toutes les petites colimites ;
  \item les colimites sont universelles (stables par tiré en arrière) ;
  \item les relations d'équivalence sont effectives ;
  \item les sommes amalgamées sont disjointes ;
  \item $\mathcal{E}$ possède une famille génératrice (petit ensemble).
  \end{enumerate}
\end{enumerate}
\end{theoreme}

\begin{proof}[Idée de preuve]
$(a) \Rightarrow (b)$ : on vérifie chaque condition pour les faisceaux.
$(b) \Rightarrow (a)$ : on prend pour $\mathcal{C}$ une petite
sous-catégorie génératrice et on munit $\mathcal{C}$ de la topologie
canonique.  Le plongement de Yoneda suivi de la faisceautisation
donne une équivalence $\mathcal{E} \simeq \mathrm{Sh}(\mathcal{C}, J)$.
\end{proof}

\index{logique interne|)}

%% ============================================================
\section{Exercices}\label{sec:exo-topos}
%% ============================================================

\begin{exercice}\label{exo:crible}
\index{crible!exercice}
Soit $\mathcal{C}$ une petite catégorie et $C \in \Ob(\mathcal{C})$.
Montrer que les cribles sur $C$ sont en bijection avec les
sous-foncteurs du foncteur représentable $\Hom(-, C)$.
\end{exercice}

\begin{exercice}\label{exo:topologie-triviale}
\index{topologie grossière!exercice}
Vérifier que la topologie grossière (seul le crible maximal est
couvrant) satisfait les axiomes (T1)--(T3).  Montrer que tout
préfaisceau est un faisceau pour cette topologie.
\end{exercice}

\begin{exercice}\label{exo:omega-Sh}
\index{classificateur de sous-objets!dans Sh(X)}
Soit $X$ un espace topologique.  Montrer que dans $\mathrm{Sh}(X)$,
le classificateur de sous-objets est le faisceau
$\Omega(U) = \{V \subseteq U \mid V \text{ ouvert}\}$
pour tout ouvert $U \subseteq X$.
En déduire que $\mathrm{Sh}(X)$ n'est en général pas booléen.
\end{exercice}

\begin{exercice}\label{exo:G-set-topos}
\index{topos!classifiant!exercice}
Soit $G$ un groupe.  Montrer que la catégorie $G\text{-}\Set$ des
$G$-ensembles (à gauche) est un topos de Grothendieck.
Décrire explicitement le classificateur $\Omega$.
\end{exercice}

\begin{exercice}[Morphismes géométriques et points]\label{exo:points-topos}
\index{point d'un topos}
Montrer qu'un morphisme géométrique $p \colon \Set \to \mathcal{E}$
(appelé \textbf{point} du topos $\mathcal{E}$) est déterminé par
un foncteur exact à gauche et qui commute aux petites colimites
$p^* \colon \mathcal{E} \to \Set$.
\end{exercice}

\begin{exercice}\label{exo:giraud-conditions}
Montrer que dans $\Set$ :
(a) les colimites sont universelles,
(b) les relations d'équivalence sont effectives,
(c) les sommes sont disjointes.
\end{exercice}

\index{topos|)}


%%%%%%%%%%%%%%%%%%%%%%%%%%%%%%%%%%%%%%%%%%%%%%%%%%%%%%%%%%%%
%%      CHAPITRE 11 : THÉORIE DES ∞-CATÉGORIES            %%
%%%%%%%%%%%%%%%%%%%%%%%%%%%%%%%%%%%%%%%%%%%%%%%%%%%%%%%%%%%%

\chapter{Théorie des $\infty$-catégories --- Aperçu}\label{ch:infty-cat}
\minitoc

\vspace{1em}

La théorie des catégories classique traite des objets et des morphismes
entre eux.  Mais dans de nombreuses situations (topologie algébrique,
algèbre homotopique, géométrie algébrique dérivée), il est naturel
de considérer des morphismes entre morphismes, des morphismes entre
ceux-ci, et ainsi de suite.  Les \textbf{$\infty$-catégories}
(ou plus précisément les $(\infty,1)$-catégories) formalisent cette
idée : ce sont des catégories avec des morphismes à tous les niveaux,
où les morphismes de niveau $\geq 2$ sont tous inversibles (à homotopie
près).  Ce chapitre final offre un aperçu des fondements, sans
prétendre à l'exhaustivité.

\index{infty-catégorie@$\infty$-catégorie|(}

%% ============================================================
\section{Motivation}\label{sec:infty-motivation}
%% ============================================================

Les catégories ordinaires ne capturent pas toute l'information
homotopique.  Considérons les exemples suivants :

\begin{enumerate}[label=(\roman*)]
\item \textbf{Espaces topologiques.}\index{type d'homotopie}
      La catégorie homotopique $\mathsf{Ho}(\Top)$ (où l'on a inversé
      les équivalences d'homotopie faible) perd de l'information :
      les espaces de morphismes ne sont plus des ensembles mais des
      \emph{espaces}.

\item \textbf{Complexes de chaînes.}\index{catégorie dérivée}
      La catégorie dérivée $D(\mathcal{A})$ d'une catégorie abélienne
      est obtenue en inversant les quasi-isomorphismes, mais cette
      localisation détruit des structures (les triangles distingués
      ne forment pas une bonne notion de limites).

\item \textbf{Champs.}\index{champ (stack)}
      Un champ sur un site est un \og faisceau de catégories\fg{}, ce
      qui demande naturellement des $2$-morphismes et des cohérences
      supérieures.
\end{enumerate}

\begin{remarque}\label{rem:modeles-infty}
\index{modèles d'$\infty$-catégories}
Il existe plusieurs modèles équivalents pour les $(\infty,1)$-catégories :
\begin{enumerate}[label=(\alph*)]
\item les quasi-catégories (Boardman--Vogt, Joyal, Lurie) ;
\item les catégories de Segal complètes (Rezk) ;
\item les catégories simpliciales (Dwyer--Kan) ;
\item les catégories de modèles (Quillen) --- comme présentation.
\end{enumerate}
Ces modèles sont tous reliés par des équivalences de Quillen ou
des équivalences catégoriques appropriées.
\end{remarque}

%% ============================================================
\section{Quasi-catégories}\label{sec:quasi-categories}
%% ============================================================

\index{quasi-catégorie|(}

\begin{notation}\label{not:simplicial}
\index{catégorie simpliciale!Delta@$\Delta$}
\index{ensemble simplicial}
On note $\Delta$ la catégorie simpliciale dont les objets sont les
ensembles ordonnés finis $[n] = \{0, 1, \ldots, n\}$ pour $n \geq 0$,
et les morphismes sont les applications croissantes.
Un \textbf{ensemble simplicial} est un foncteur
$X \colon \Delta^\op \to \Set$.
On note $\mathbf{sSet} = \Fun(\Delta^\op, \Set)$ la catégorie
des ensembles simpliciaux.
\end{notation}

\begin{definition}[Corne]\label{def:corne}
\index{corne}
\index{Lambda@$\Lambda^n_k$ (corne)}
Pour $0 \leq k \leq n$, la \textbf{corne}
$\Lambda^n_k \subset \Delta^n$ est le sous-ensemble simplicial
obtenu en retirant la face intérieure et la $k$-ème face de
$\Delta^n$.  La corne est dite :
\begin{enumerate}[label=(\roman*)]
\item \textbf{intérieure} si $0 < k < n$ ;
\item \textbf{extérieure} si $k = 0$ ou $k = n$.
\end{enumerate}
\end{definition}

\begin{definition}[Quasi-catégorie]\label{def:quasi-categorie}
\index{quasi-catégorie!définition}
\index{complexe de Kan faible}
\index{Joyal}
Un ensemble simplicial $X$ est une \textbf{quasi-catégorie}
(ou \textbf{$\infty$-catégorie} au sens de Joyal--Lurie) s'il
satisfait la \textbf{condition de relèvement de Kan interne} :
pour tout $n \geq 2$ et tout $0 < k < n$, toute application
$\Lambda^n_k \to X$ admet une extension à $\Delta^n$ :
\[
\begin{tikzcd}
\Lambda^n_k \arrow[r] \arrow[d, hook] & X \\
\Delta^n \arrow[ur, dashed] &
\end{tikzcd}
\]
Contrairement à un complexe de Kan (qui exige le relèvement pour
\emph{toutes} les cornes, y compris extérieures), une quasi-catégorie
ne requiert le relèvement que pour les cornes intérieures.
\end{definition}

\begin{exemple}[Quasi-catégories fondamentales]\label{ex:quasi-categories}
\index{quasi-catégorie!exemples}
\begin{enumerate}[label=(\roman*)]
\item \textbf{Nerf d'une catégorie.}\index{nerf d'une catégorie}
      Si $\mathcal{C}$ est une catégorie ordinaire, son nerf
      $N(\mathcal{C})$ est une quasi-catégorie.
      Les $0$-simplexes sont les objets, les $1$-simplexes sont les
      morphismes, et les $n$-simplexes ($n \geq 2$) sont les suites
      composables de $n$ morphismes.
      Le relèvement interne est \emph{unique} car la composition dans
      $\mathcal{C}$ est strictement définie.

\item \textbf{Complexes de Kan.}\index{complexe de Kan}
      Tout complexe de Kan est une quasi-catégorie.
      Les complexes de Kan correspondent aux $\infty$-groupoïdes
      (tous les morphismes sont inversibles à homotopie près).
      L'ensemble simplicial singulier $\mathrm{Sing}(X)$ d'un espace
      topologique $X$ est un complexe de Kan.

\item \textbf{Nerf de Dold--Kan.}
      \index{nerf de Dold--Kan}
      Le nerf de la catégorie dérivée d'une catégorie abélienne
      donne une quasi-catégorie qui encode l'information homotopique
      perdue par la localisation.
\end{enumerate}
\end{exemple}

\begin{proposition}[Catégorie homotopique]\label{prop:cat-homotopique}
\index{catégorie homotopique!d'une quasi-catégorie}
Soit $X$ une quasi-catégorie.  On définit sa \textbf{catégorie
homotopique} $\mathrm{h}X$ comme la catégorie ordinaire dont :
\begin{enumerate}[label=(\roman*)]
\item les objets sont les $0$-simplexes de $X$ ;
\item les morphismes de $x$ à $y$ sont les classes d'homotopie de
      $1$-simplexes de $x$ à $y$ ;
\item la composition est induite par le relèvement des cornes
      $\Lambda^2_1$.
\end{enumerate}
La catégorie $\mathrm{h}X$ est bien définie et
$\mathrm{h}(N(\mathcal{C})) \simeq \mathcal{C}$ pour toute catégorie
ordinaire $\mathcal{C}$.
\end{proposition}

\index{quasi-catégorie|)}

%% ============================================================
\section{$(\infty,1)$-Catégories}\label{sec:infty-1-categories}
%% ============================================================

\index{(infty,1)-catégorie@$(\infty,1)$-catégorie|(}

\begin{definition}[$(\infty,1)$-catégorie (informelle)]
\label{def:infty-1-cat}
\index{(infty,1)-catégorie@$(\infty,1)$-catégorie!définition}
Une \textbf{$(\infty,1)$-catégorie} est une structure possédant :
\begin{enumerate}[label=(\roman*)]
\item des objets ;
\item des $1$-morphismes entre objets ;
\item des $2$-morphismes entre $1$-morphismes ;
\item des $n$-morphismes entre $(n-1)$-morphismes, pour tout $n \geq 2$ ;
\end{enumerate}
où tous les $k$-morphismes pour $k \geq 2$ sont inversibles
(à homotopie près).
\end{definition}

\begin{remarque}\label{rem:infty-vs-quasi}
Le modèle standard pour les $(\infty,1)$-catégories est celui
des quasi-catégories.  Dans la suite, on utilise les termes
\og $\infty$-catégorie\fg{} et \og quasi-catégorie\fg{} de manière
interchangeable.
\end{remarque}

\begin{definition}[Espaces de morphismes]\label{def:mapping-space}
\index{espace de morphismes}
\index{mapping space}
Soit $\mathcal{C}$ une $\infty$-catégorie et $x, y \in \mathcal{C}$.
L'\textbf{espace de morphismes} $\mathrm{Map}_\mathcal{C}(x, y)$
est un espace de Kan (un complexe de Kan) défini par
\[
  \mathrm{Map}_\mathcal{C}(x, y)_n
  = \{f \colon \Delta^{n+1} \to \mathcal{C} \mid
      f|_{\Delta^{\{0\}}} = x,\;
      f|_{\Delta^{\{1,\ldots,n+1\}}} \text{ dégénéré sur } y\}.
\]
Les composantes connexes $\pi_0(\mathrm{Map}_\mathcal{C}(x,y))$
redonnent l'ensemble des morphismes dans $\mathrm{h}\mathcal{C}$.
\end{definition}

\begin{exemple}[$\infty$-catégories fondamentales]\label{ex:infty-cats}
\index{(infty,1)-catégorie@$(\infty,1)$-catégorie!exemples}
\begin{enumerate}[label=(\roman*)]
\item \textbf{$\mathcal{S}$ --- l'$\infty$-catégorie des espaces.}
      \index{S@$\mathcal{S}$ ($\infty$-catégorie des espaces)}
      Les objets sont les complexes de Kan (ou les espaces
      topologiques, à équivalence faible près).  Les espaces de
      morphismes sont les espaces de fonctions.

\item \textbf{$\mathcal{C}\mathrm{at}_\infty$ --- l'$\infty$-catégorie
      des $\infty$-catégories.}
      \index{Cat infty@$\mathcal{C}\mathrm{at}_\infty$}
      C'est l'analogue supérieur de $\Cat$.

\item \textbf{$\mathcal{D}(\mathcal{A})$ --- catégorie dérivée
      $\infty$-catégorique.}
      \index{catégorie dérivée!infty@$\infty$-catégorique}
      Pour une catégorie abélienne $\mathcal{A}$, on obtient une
      $\infty$-catégorie stable qui raffine la catégorie dérivée
      triangulée classique.
\end{enumerate}
\end{exemple}

%% ============================================================
\section{Yoneda $\infty$-catégorique}\label{sec:yoneda-infty}
%% ============================================================

\begin{theoreme}[Lemme de Yoneda $\infty$-catégorique]
\label{thm:yoneda-infty}
\index{lemme de Yoneda!infty-catégorique@$\infty$-catégorique}
\index{Lurie}
Soit $\mathcal{C}$ une petite $\infty$-catégorie.  Le plongement
de Yoneda
\[
  \mathbf{y} \colon \mathcal{C} \longrightarrow
  \Fun(\mathcal{C}^\op, \mathcal{S})
\]
est pleinement fidèle : pour tous $x, y \in \mathcal{C}$, l'application
\[
  \mathrm{Map}_\mathcal{C}(x, y)
  \xrightarrow{\;\sim\;}
  \mathrm{Map}_{\Fun(\mathcal{C}^\op, \mathcal{S})}(\mathbf{y}(x), \mathbf{y}(y))
\]
est une équivalence d'espaces.
\end{theoreme}

\begin{remarque}\label{rem:yoneda-infty-subtil}
La preuve du lemme de Yoneda $\infty$-catégorique est considérablement
plus subtile que dans le cas classique.
Lurie en donne une démonstration dans \emph{Higher Topos Theory}
(Proposition 5.1.3.1) en utilisant les joints d'ensembles
simpliciaux et les extensions de Kan à droite.
\end{remarque}

%% ============================================================
\section{Adjonctions $\infty$-catégoriques}\label{sec:adjonctions-infty}
%% ============================================================

\begin{definition}[Adjonction entre $\infty$-catégories]
\label{def:adjonction-infty}
\index{adjonction!infty-catégorique@$\infty$-catégorique}
Soient $\mathcal{C}$ et $\mathcal{D}$ deux $\infty$-catégories.
Un foncteur $F \colon \mathcal{C} \to \mathcal{D}$ est
\textbf{adjoint à gauche} d'un foncteur
$G \colon \mathcal{D} \to \mathcal{C}$ s'il existe une
équivalence d'espaces
\[
  \mathrm{Map}_\mathcal{D}(F(x), y)
  \simeq
  \mathrm{Map}_\mathcal{C}(x, G(y))
\]
naturelle en $x \in \mathcal{C}$ et $y \in \mathcal{D}$.
\end{definition}

\begin{proposition}[Caractérisation par unité/coünité]
\label{prop:adjonction-infty-unite}
\index{adjonction!unité et coünité ($\infty$)}
L'adjonction $F \dashv G$ dans le cadre $\infty$-catégorique est
équivalente à l'existence de transformations naturelles
$\eta \colon \id_\mathcal{C} \to G \circ F$ et
$\varepsilon \colon F \circ G \to \id_\mathcal{D}$
satisfaisant les identités triangulaires à homotopie (cohérente) près.
\end{proposition}

%% ============================================================
\section{Topos supérieurs}\label{sec:topos-superieurs}
%% ============================================================

\index{topos!supérieur|(}

\begin{definition}[Topos supérieur (d'après Lurie)]
\label{def:topos-superieur}
\index{topos!supérieur!définition}
\index{Lurie!topos supérieur}
Un \textbf{$\infty$-topos} (ou topos supérieur) est une
$\infty$-catégorie $\mathcal{X}$ qui est une localisation exacte
à gauche accessible d'une $\infty$-catégorie de préfaisceaux :
il existe une petite $\infty$-catégorie $\mathcal{C}$ et une
localisation
\[
  L \colon \Fun(\mathcal{C}^\op, \mathcal{S})
  \rightleftarrows \mathcal{X} \cocolon \iota
\]
avec $L \dashv \iota$ et $L$ exact à gauche.
\end{definition}

\begin{theoreme}[Giraud $\infty$-catégorique (Lurie)]
\label{thm:giraud-infty}
\index{théorème de Giraud!infty-catégorique@$\infty$-catégorique}
Soit $\mathcal{X}$ une $\infty$-catégorie présentable.
Les conditions suivantes sont équivalentes :
\begin{enumerate}[label=(\alph*)]
\item $\mathcal{X}$ est un $\infty$-topos ;
\item les colimites dans $\mathcal{X}$ sont universelles, et les
      objets de groupoïde sont effectifs ;
\item $\mathcal{X}$ vérifie la descente : pour tout diagramme
      simplicial $U_\bullet \to X$, le foncteur
      $\mathcal{X}_{/X} \to \lim_n \mathcal{X}_{/U_n}$ est une
      équivalence.
\end{enumerate}
\end{theoreme}

\begin{exemple}[$\infty$-topos fondamentaux]\label{ex:infty-topos}
\index{topos!supérieur!exemples}
\begin{enumerate}[label=(\roman*)]
\item $\mathcal{S}$, l'$\infty$-catégorie des espaces, est le
      $\infty$-topos terminal (analogue de $\Set$).
\item Pour un espace topologique $X$, la catégorie
      $\mathrm{Sh}_\infty(X)$ des faisceaux d'espaces sur $X$ est
      un $\infty$-topos.
\item Le $\infty$-topos étale d'un schéma intervient en géométrie
      algébrique dérivée.
      \index{géométrie algébrique dérivée}
\end{enumerate}
\end{exemple}

\begin{remarque}[Perspectives]\label{rem:perspectives}
La théorie des $\infty$-catégories et des $\infty$-topos est
aujourd'hui un domaine très actif.  Parmi les directions de
recherche :
\begin{enumerate}[label=(\roman*)]
\item la \textbf{géométrie algébrique dérivée} (Lurie, Toën--Vezzosi) ;
\item la \textbf{topologie algébrique chromatique}
      \index{homotopie chromatique}
      (spectres, localisation chromatique) ;
\item les \textbf{théories des champs topologiques}
      (TQFT, théorie de Cobordisme étendu,
      hypothèse du cobordisme de Baez--Dolan--Lurie) ;
      \index{hypothèse du cobordisme}
\item la \textbf{logique homotopique} et la théorie des types
      homotopiques (HoTT/UF) ;
      \index{HoTT (théorie des types homotopiques)}
\item les \textbf{$(\infty,n)$-catégories} pour $n > 1$ et les
      structures de dimension supérieure.
\end{enumerate}
\end{remarque}

\index{topos!supérieur|)}

%% ============================================================
\section{Exercices}\label{sec:exo-infty}
%% ============================================================

\begin{exercice}\label{exo:nerf-quasi-cat}
\index{nerf d'une catégorie!exercice}
Soit $\mathcal{C}$ une catégorie ordinaire.  Montrer que le nerf
$N(\mathcal{C})$ est une quasi-catégorie, et que les relèvements
de cornes intérieures sont \emph{uniques}.
\end{exercice}

\begin{exercice}\label{exo:kan-infty-groupoide}
\index{complexe de Kan!et $\infty$-groupoïde}
Soit $X$ un complexe de Kan.  Montrer que la catégorie homotopique
$\mathrm{h}X$ est un groupoïde (tout morphisme est un isomorphisme).
\end{exercice}

\begin{exercice}\label{exo:mapping-pi0}
Soit $\mathcal{C}$ une quasi-catégorie et $x, y \in \mathcal{C}$.
Montrer que $\pi_0(\mathrm{Map}_\mathcal{C}(x,y)) \cong
\Hom_{\mathrm{h}\mathcal{C}}(x, y)$.
\end{exercice}

\begin{exercice}\label{exo:infty-topos-S}
\index{S@$\mathcal{S}$!comme $\infty$-topos}
Montrer que l'$\infty$-catégorie $\mathcal{S}$ des espaces est le
$\infty$-topos final : pour tout $\infty$-topos $\mathcal{X}$, il
existe un morphisme géométrique essentiellement unique
$\mathcal{X} \to \mathcal{S}$.
\end{exercice}

\begin{exercice}[Nerf et catégorie homotopique]\label{exo:nerf-h}
Montrer que pour toute catégorie ordinaire $\mathcal{C}$,
le foncteur catégorie homotopique du nerf redonne la catégorie
d'origine : $\mathrm{h}(N(\mathcal{C})) \cong \mathcal{C}$.
\end{exercice}

\begin{exercice}\label{exo:sing-kan}
\index{ensemble simplicial singulier}
Soit $X$ un espace topologique.  Montrer que l'ensemble simplicial
singulier $\mathrm{Sing}(X)$ est un complexe de Kan.
\textit{Indication} : utiliser le fait que les cornes topologiques
sont des rétractes de déformation des simplexes standard.
\end{exercice}

\index{infty-catégorie@$\infty$-catégorie|)}
\index{(infty,1)-catégorie@$(\infty,1)$-catégorie|)}


%%%%%%%%%%%%%%%%%%%%%%%%%%%%%%%%%%%%%%%%%%%%%%%%%%%%%%%%%%%%
%%                     APPENDICE                           %%
%%%%%%%%%%%%%%%%%%%%%%%%%%%%%%%%%%%%%%%%%%%%%%%%%%%%%%%%%%%%

\appendix

\chapter{Rappels d'algèbre et de topologie}\label{app:rappels}
\minitoc

\vspace{1em}

Cet appendice rassemble sans démonstration les résultats d'algèbre
et de topologie utilisés dans le cours.

%% ============================================================
\section{Algèbre}\label{sec:rappels-algebre}
%% ============================================================

\index{rappels d'algèbre|(}

\begin{definition}[Groupe, anneau, module]\label{def:rappel-structures}
\index{groupe!rappel}
\index{anneau!rappel}
\index{module!rappel}
\begin{enumerate}[label=(\roman*)]
\item Un \textbf{groupe} $(G, \cdot, e)$ est un ensemble muni d'une
      loi associative, d'un élément neutre et d'inverses.

\item Un \textbf{anneau} $(R, +, \cdot, 0, 1)$ est un groupe abélien
      $(R, +, 0)$ muni d'une seconde loi associative avec unité $1$,
      distributive par rapport à l'addition.

\item Un \textbf{$R$-module à gauche} est un groupe abélien $(M, +)$
      muni d'une action $R \times M \to M$ compatible avec les
      structures de $R$ et de $M$.
\end{enumerate}
\end{definition}

\begin{proposition}[Produit tensoriel]\label{prop:rappel-tensoriel}
\index{produit tensoriel!rappel}
Soient $R$ un anneau commutatif, $M$ un $R$-module à droite et $N$
un $R$-module à gauche.  Le \textbf{produit tensoriel}
$M \otimes_R N$ est le quotient du $\Z$-module libre sur
$M \times N$ par les relations de bilinéarité et de $R$-équilibre :
\[
  (m \cdot r) \otimes n = m \otimes (r \cdot n).
\]
Il satisfait la propriété universelle :
$\Hom_\Z(M \otimes_R N, P) \cong \mathrm{Bil}_R(M \times N, P)$.
\end{proposition}

\begin{definition}[Suite exacte]\label{def:rappel-suite-exacte}
\index{suite exacte!rappel}
Une suite de morphismes de $R$-modules
$\cdots \to M_{n+1} \xrightarrow{f_{n+1}} M_n
\xrightarrow{f_n} M_{n-1} \to \cdots$
est \textbf{exacte} en $M_n$ si $\ker f_n = \mathrm{im}\, f_{n+1}$.
Une \textbf{suite exacte courte} est une suite exacte
$0 \to A \xrightarrow{f} B \xrightarrow{g} C \to 0$.
\end{definition}

\begin{definition}[Algèbre libre, quotient]\label{def:rappel-algebre-libre}
\index{algèbre libre!rappel}
\index{quotient!rappel}
Soit $R$ un anneau commutatif.
\begin{enumerate}[label=(\roman*)]
\item Le $R$-module libre sur un ensemble $S$ est
      $R^{(S)} = \bigoplus_{s \in S} R$.
\item Pour un sous-module $N \subseteq M$, le quotient $M/N$ est
      un $R$-module avec la surjection canonique $\pi \colon M \to M/N$.
\end{enumerate}
\end{definition}

\index{rappels d'algèbre|)}

%% ============================================================
\section{Topologie}\label{sec:rappels-topologie}
%% ============================================================

\index{rappels de topologie|(}

\begin{definition}[Espace topologique]\label{def:rappel-topologique}
\index{espace topologique!rappel}
Un \textbf{espace topologique} est un couple $(X, \tau)$ où $\tau$
est une collection de sous-ensembles de $X$ (les \textbf{ouverts})
stable par unions quelconques et intersections finies, contenant
$\varnothing$ et $X$.
\end{definition}

\begin{definition}[Application continue, homéomorphisme]
\label{def:rappel-continue}
\index{application continue!rappel}
\index{homéomorphisme!rappel}
Une application $f \colon X \to Y$ entre espaces topologiques est
\textbf{continue} si $f^{-1}(V)$ est ouvert dans $X$ pour tout
ouvert $V$ de $Y$.  Un \textbf{homéomorphisme} est une bijection
continue dont l'inverse est continue.
\end{definition}

\begin{definition}[Compacité, connexité]\label{def:rappel-compact-connexe}
\index{compacité!rappel}
\index{connexité!rappel}
\begin{enumerate}[label=(\roman*)]
\item $X$ est \textbf{compact} si de tout recouvrement ouvert on peut
      extraire un sous-recouvrement fini.
\item $X$ est \textbf{connexe} s'il ne peut pas s'écrire comme
      union de deux ouverts disjoints non vides.
\end{enumerate}
\end{definition}

\begin{definition}[Préfaisceau et faisceau sur un espace]
\label{def:rappel-faisceau}
\index{faisceau!rappel}
\index{préfaisceau!rappel}
Soit $X$ un espace topologique.
\begin{enumerate}[label=(\roman*)]
\item Un \textbf{préfaisceau} (d'ensembles) sur $X$ est un foncteur
      $F \colon \mathsf{Ouv}(X)^\op \to \Set$.
\item Un préfaisceau $F$ est un \textbf{faisceau} s'il satisfait la
      condition de recollement : pour tout recouvrement ouvert
      $\{U_i\}$ de $U$, le diagramme
      \[
        F(U) \to \prod_i F(U_i) \rightrightarrows
        \prod_{i,j} F(U_i \cap U_j)
      \]
      est un égalisateur.
\end{enumerate}
\end{definition}

\begin{definition}[Homotopie, équivalence d'homotopie]
\label{def:rappel-homotopie}
\index{homotopie!rappel}
\index{équivalence d'homotopie!rappel}
\begin{enumerate}[label=(\roman*)]
\item Deux applications continues $f, g \colon X \to Y$ sont
      \textbf{homotopes} s'il existe une application continue
      $H \colon X \times [0,1] \to Y$ avec $H(-, 0) = f$ et
      $H(-, 1) = g$.
\item $f \colon X \to Y$ est une \textbf{équivalence d'homotopie}
      s'il existe $g \colon Y \to X$ avec $g \circ f \simeq \id_X$
      et $f \circ g \simeq \id_Y$.
\end{enumerate}
\end{definition}

\begin{definition}[Groupe fondamental]\label{def:rappel-pi1}
\index{groupe fondamental!rappel}
Pour un espace topologique pointé $(X, x_0)$, le \textbf{groupe
fondamental} $\pi_1(X, x_0)$ est le groupe des classes d'homotopie
de lacets basés en $x_0$.
\end{definition}

\index{rappels de topologie|)}

%% ============================================================
%%                       INDEX
%% ============================================================

\printindex

\end{document}
